\chapter{Moral Graph}
\label{ch-moral-g}

We define  the query $Q=\rvx\perp\rvz\mid \rvy?$
to be true if

\beq
P(x| z, y)=P(x|y)
\eeq
In Chapter \ref{ch-dsep},
we gave necessary and sufficient
graphical conditions on a bnet
for the query $Q=\rvx\perp\rvz\mid \rvy$ to be true.
When those conditions, called d-separation conditions, are
satisfied, we will say $\rvx\perp_{bnet} \rvz\mid \rvy$ is true. Thus

\beq
\rvx\perp \rvz\mid \rvy
\iff
\rvx\perp_{bnet} \rvz\mid \rvy
\eeq

In this Chapter, we will show that given
a bnet $\calb$ and a query $Q$, one can construct an
undirected graph called the {\bf ancestral
moral (AM) graph},
denoted by $AM(\calb, Q)$.
For an AM graph, $\rvx\perp_{AM}\rvz\mid \rvy$
iff all (undirected) paths from
$\rvx$ to $\rvz$ are blocked by a node in $\rvy$. In general, $\rvy$ can be not just a single node, but a set of nodes (i.e., multi-node) or the empty set.

In this chapter we will also
show that

\beq
\boxed{\rvx\perp_{bnet} \rvz\mid \rvy
\iff
\rvx\perp_{AM} \rvz\mid \rvy}
\label{eq-perp-bnet-AM}
\eeq
Recall that the d-separation conditions for
a bnet handle
colliders, mediators and  forks differently.
A nice feature of the d-separation conditions for an AM graph is that
they are the same for colliders, mediators and  forks, as one might
expect, since there is no difference
between those 3 in an undirected graph.




\begin{figure}[h!]
$$
\begin{array}{ccc}
\xymatrix{
\rva\ar[r]\ar[d]
&{\rvb}\ar[d]
\\
{\rvc}\ar[r]\ar[d]
&{\rvd}\ar[d]
\\
\rvf
&\rve\ar[l]
}
&&
\xymatrix{
\rva\ar@{-}[r]\ar@{-}[d]
&{\rvb}\ar@{-}[d]\ar@{-}[dl]
\\
{\rvc}\ar@{-}[r]\ar@{-}[d]\ar@{-}[dr]
&{\rvd}\ar@{-}[d]
\\
\rvf
&\rve\ar@{-}[l]
}
\\
(a)
&& (b)
\end{array}
$$
\caption{Example of construction of a moral graph $(b)$ from a bnet $(a)$.}
\label{fig-moral-g}
\end{figure}
Fig.\ref{fig-moral-g}
shows an example of how to construct a
{\bf moral graph} $Moral(\calb)$ from a bnet $\calb$.
To do so,
\begin{enumerate}
\item Add an undirected line (edge) between all
pairs of \qt{unwed} (i.e., not connected by an arrow already) parents of all colliders. If the collider
has 2 parents, then one adds $2!/2=1$ edges.
If the collider has 3 parents, then one adds $3!/2=3$ edges which form a clique (i.e.,
fully connected subgraph). And so forth
for more than 3 parents.
\item Remove arrow directions
\end{enumerate}


\begin{figure}
$$
\begin{array}{cccccc}
\xymatrix@C=1pc{
\rva\ar[r]\ar[d]
&\Circle{\rvb}\ar[d]
\\
\Circle{\rvc}\ar[r]\ar[d]
&\Circle{\rvd}\ar[d]
\\
\rvf
&\rve\ar[l]
}
&
\xymatrix@C=1pc{
\Circle{\rva}\ar[r]\ar[d]
&\Circle{\rvb}\ar[d]
\\
\Circle{\rvc}\ar[r]\ar[d]
&\Circle{\rvd}\ar[d]
\\
\rvf
&\rve\ar[l]
}
&
\xymatrix@C=1pc{
\Circle{\rva}\ar@[red][r]\ar@[red][d]
&\Circle{\rvb}\ar@[red][d]
\\
\Circle{\rvc}\ar@[red][r]\ar[d]
&\Circle{\rvd}\ar[d]
\\
\rvf
&\rve\ar[l]
}
&
\xymatrix@C=1pc{
\Circle{\rva}\ar@[red][r]\ar@[red][d]
&\Circle{\rvb}\ar@[red][d]
\\
\Circle{\rvc}\ar@[red][r]
&\Circle{\rvd}
}
&
\xymatrix@C=1pc{
\Circle{\rva}\ar@[red][r]\ar@[red][d]
&\Circle{\rvb}\ar@[red][d]
\ar@[red]@{-}[dl]
\\
\Circle{\rvc}\ar@[red][r]
&\Circle{\rvd}
}
&
\xymatrix@C=1pc{
\Circle{\rva}\ar@{-}@[red][r]\ar@{-}@[red][d]
&\Circle{\rvb}\ar@{-}@[red][d]
\ar@{-}@[red][dl]
\\
\Circle{\rvc}\ar@{-}@[red][r]
&\Circle{\rvd}
}
\\
\\
(1)
&(2)
&(3)
&(4)
&(5)
&(6)
\end{array}
$$
\caption{Construction of ancestral moral graph
for query $\rvc\perp \rvb\mid \rvd?$}
\label{fig-anc-moral-c-b-d}
\end{figure}

Fig.\ref{fig-anc-moral-c-b-d}
show an example of how to construct the Ancestral Moral (AM) graph
$AM(\calb, Q)$
of a given bnet $\calb$ for the query
$Q=\rvc\perp \rvb\mid \rvd?$.
In that figure
we perform the following steps

\begin{enumerate}
\item Circle the 3 nodes
$\rvb, \rvc, \rvd$  involved
in the query $Q=\rvc\perp \rvb\mid \rvd?$.
The set $S_Q=\{\rvb, \rvc, \rvd\}$
is called the {\bf scope} of the query.

\item Add to the scope of the  query
all ancestors of $S_Q$, including $S_Q$ itself.
This set of nodes will be denoted by $An(S_Q)$.

\item Mark in red all arrows, and only the arrows between the
nodes in $An(S_Q)$.

\item Remove non-red arrows

\item Add an undirected line between all
pairs of \qt{unwed} (i.e., not connected by an arrow already) parents of all colliders.

\item Remove the direction of all arrows.

\end{enumerate}

We will not prove rigorously the
equivalence Eq.(\ref{eq-perp-bnet-AM})
between the bnet and the AM
d-separation
conditions.
Instead, we give the
following examples to
convince the reader
that it's true.
Figs.\ref{fig-moral-fork}/
\ref{fig-moral-mediator}/
\ref{fig-moral-collider}
give some
examples of application of both, the
bnet and AM d-separation conditions, to  fork/mediator/collider nodes.
\newpage
\begin{figure}[h!]
$$
\begin{tabular}{c|c|c|c|c|c|}
{bnet}
&{query}
&\begin{tabular}{c}
{answer}
\\
{to query}
\\
(bnet
\\
d-sep)
\end{tabular}
&
{scope}
&
\begin{tabular}{c}
{ancestral}
\\
{moral}
\\
{graph}
\end{tabular}
&\begin{tabular}{c}
{answer}
\\
{to query}
\\
(AM
\\
d-sep)
\end{tabular}
\\ \hline\hline %row 1
\xymatrix@C=1pc{
\rvx
&*++[o][F*:yellow]{\rvy}\ar[l]\ar[r]
&\rvz
}
& $\rvx\perp \rvz\mid \rvy?$
&\text{Yes}
&$\{\rvx, \rvy,\rvz\}$
& \xymatrix@C=1pc{\rvx
	&*++[o][F*:yellow]{\rvy}\ar@{-}[l]\ar@{-}[r]
	&\rvz}
& \text{Yes}
\\
\hline %row 2
\xymatrix@C=1pc{
\rvx
&\rvy\ar[l]\ar[r]
&\rvz
}
& $\rvx\perp \rvz?$
&\text{No}
&$\{\rvx, \rvz\}$
& \xymatrix@C=1pc{\rvx
	&\rvy\ar@{-}[l]\ar@{-}[r]
	&\rvz}
& \text{No}
\\
\hline %row 3
\begin{tabular}{c}
\xymatrix@C=1pc{
\rvx\ar@/_1pc/[rr]
&*++[o][F*:yellow]{\rvy}\ar[l]\ar[r]
&\rvz
}
\\
\\
\end{tabular}
& $\rvx\perp \rvz\mid \rvy?$
&\text{No}
&$\{\rvx, \rvy,\rvz\}$
& \xymatrix@C=1pc{\rvx\ar@/_1pc/@{-}[rr]
	&*++[o][F*:yellow]{\rvy}\ar@{-}[l]\ar@{-}[r]
	&\rvz}
& \text{No}
\\
\hline %row 4
\begin{tabular}{c}
\xymatrix@C=1pc{
\rvx\ar@/_1pc/[rr]
&\rvy\ar[l]\ar[r]
&\rvz
}
\\
\\
\end{tabular}
& $\rvx\perp \rvz?$
&\text{No}
&$\{\rvx, \rvz\}$
& \xymatrix@C=1pc{\rvx\ar@/_1pc/@{-}[rr]
	&\rvy\ar@{-}[l]\ar@{-}[r]
	&\rvz}
& \text{No}
\\ \hline
\end{tabular}
$$
\caption{{Examples of application of both, the  bnet and AM d-separation conditions, to a fork node.}}
\label{fig-moral-fork}
\end{figure}


\begin{figure}[h!]
$$e
\begin{tabular}{c|c|c|c|c|c|}
{bnet}
&{query}
&\begin{tabular}{c}
{answer}
\\
{to query}
\\
(bnet
\\
d-sep)
\end{tabular}
&
{scope}
&
\begin{tabular}{c}
{ancestral}
\\
{moral}
\\
{graph}
\end{tabular}
&\begin{tabular}{c}
{answer}
\\
{to query}
\\
(AM
\\
d-sep)
\end{tabular}
\\ \hline\hline %row 1
\xymatrix@C=1pc{
\rvx\ar[r]
&*++[o][F*:yellow]{\rvy}\ar[r]
&\rvz
}
& $\rvx\perp \rvz\mid \rvy?$
&\text{Yes}
&$\{\rvx, \rvy,\rvz\}$
& \xymatrix@C=1pc{\rvx
	&*++[o][F*:yellow]{\rvy}\ar@{-}[l]\ar@{-}[r]
	&\rvz}
& \text{Yes}
\\
\hline %row 2
\xymatrix@C=1pc{
\rvx\ar[r]
&\rvy\ar[r]
&\rvz
}
& $\rvx\perp \rvz?$
&\text{No}
&$\{\rvx, \rvz\}$
& \xymatrix@C=1pc{\rvx
	&{\rvy}\ar@{-}[l]\ar@{-}[r]
	&\rvz}
& \text{No}
\\
\hline %row 3
\begin{tabular}{c}
\xymatrix@C=1pc{
\rvx\ar@/_1pc/[rr]\ar[r]
&*++[o][F*:yellow]{\rvy}\ar[r]
&\rvz
}
\\
\\
\end{tabular}
& $\rvx\perp \rvz\mid \rvy?$
&\text{No}
&$\{\rvx, \rvy,\rvz\}$
& \xymatrix@C=1pc{\rvx\ar@/_1pc/@{-}[rr]
	&*++[o][F*:yellow]{\rvy}\ar@{-}[l]\ar@{-}[r]
	&\rvz}
& \text{No}
\\
\hline %row 4
\begin{tabular}{c}
\xymatrix@C=1pc{
\rvx\ar@/_1pc/[rr]
\ar[r]
&\rvy\ar[r]
&\rvz
}
\\
\\
\end{tabular}
& $\rvx\perp \rvz?$
&\text{No}
&$\{\rvx, \rvz\}$
& \xymatrix@C=1pc{\rvx\ar@/_1pc/@{-}[rr]
	&\rvy\ar@{-}[l]\ar@{-}[r]
	&\rvz}
& \text{No}
\\ \hline
\end{tabular}
$$
\caption{Examples of application of both, the
bnet and AM
 d-separation conditions, to a mediator node.}
\label{fig-moral-mediator}
\end{figure}
\newpage
\begin{figure}[h!]
$$
\begin{tabular}[t]{l|l|l|l|l|l|}
{bnet}
&{query}
&\begin{tabular}{c}
{answer}
\\
{to query}
\\
(bnet
\\
d-sep)
\end{tabular}
&
{scope}
&
\begin{tabular}{c}
{ancestral}
\\
{moral}
\\
{graph}
\end{tabular}
&\begin{tabular}{c}
{answer}
\\
{to query}
\\
(AM
\\
d-sep)
\end{tabular}
\\ \hline\hline %row 1
\xymatrix@R=1pc@C=1pc{
&\rva\ar[d]
\\
&\rvb
\ar[dr]
\\
\rvx\ar[dr]\ar[ur]
&&\rvz\ar[dl]
\\
&*++[o][F*:yellow]{\rvy}\ar[d]
\\
\rve\ar[r]&\rvd
}
&$\rvx\perp\rvz\mid \rvy?$
&No
&$\{\rvx, \rvy, \rvz\}$
&\xymatrix@R=1pc@C=1pc{
&\rva
\ar@{-}[d]
\\
&\rvb
\\
\rvx\ar@{-}[dr]
&&\rvz\ar@{-}[dl]\ar@{-}[ul]
\ar@{-}[ll]
\\
&*++[o][F*:yellow]{\rvy}
}
&No
\\
\hline %row 2
\xymatrix@R=1pc@C=1pc{ &\rva\ar[d] \\ &\rvb\ar[dr] \\ \rvx\ar[dr] \ar[ur] &&\rvz\ar[dl]  \\ &\rvy\ar[d] \\ \rve\ar[r] &*++[o][F*:yellow]{\rvd} }
& $\rvx\perp \rvz\mid \rvd?$
&No
&$\{\rvx, \rvz, \rvd\}$
&\xymatrix@R=1pc@C=1pc{ &\rva\ar@{-}[d] \\ &\rvb \\ \rvx\ar@{-}[dr]  &&\rvz\ar@{-}[dl] \ar@{-}[ll] \ar@{-}[ul] \\ &\rvy\ar@{-}[d] \\ \rve\ar@{-}[r] &*++[o][F*:yellow]{\rvd} }
& No
\\
\hline %row 3
\xymatrix@R=1pc@C=1pc{ &\rva\ar[d] \\ &\rvb \\ \rvx\ar[dr] \ar[ur] &&\rvz\ar[dl] \ar[ul] \\ &\rvy\ar[d] \\ \rve\ar[r] &\rvd }
& $\rvx\perp \rvz?$
&Yes
&$\{\rvx, \rvz\}$
&
\xymatrix@C=1pc@R=1pc{
&
\\
&
\\
\rvx
&&\rvz
}
& Yes
\\ \hline
\end{tabular}
$$
\caption{Examples of application of
both, the
bnet and AM d-separation conditions, to a collider node.}
\label{fig-moral-collider}
\end{figure}
\hrule
AM graphs in hypergraph language

A {\bf powerset} $\calp(S)$ of a set $S$
is the collection of all subsets of $S$ including $S$ and $\eset$.
For example, $\calp(\{1,2\})= \{\eset, \{1\}, \{2\}, \{1,2\}\}$.

A {\bf hypergraph} $\calh$ is a subset of a powerset $\calp(S)$. Hence, $\calh\subset \calp(S)$. For example, $\{\eset, \{1\}, \{1,2\}\}$ is a hypergraph. A hypergraph can be specified
in proper Graph Theory format by saying that
it's a pair $(V, E)$ of vertices $V$ and edges $E$, where $E=\calh$ and $V$ are the vertices mentioned in $E$.

Powersets can be charted with lattice graphs, as shown in Chapter \ref{ch-poset}. In such
lattice graphs, each row corresponds to sets of a given order (order of a set=number of elements). Also, that order decreases by 1 as we go down the lattice graph
by 1 row. The row with sets of order $p$ is
called the {\bf $p$-section of the lattice} $\calp(S)$ (and the intersection of that  $p$-section with a hypergraph $\calh\subset \calp(S)$ is called the {\bf $p$-section
of the hypergraph} $\calh$).


For a bnet $\calb$, define
the {\bf family  of a node $\rvx$} as the set $Fam(\rvx)$ consisting of $\rvx$ and its parents.
Define the {\bf family hypergraph}
for bnet $\calb$
with node set $nds(\calb)$ as

\beq
\calh_{fam}(\calb)=
\{Fam(\rvx) :\rvx\in nds(\calb)\}
\eeq
For example

\beq
\text{If}\quad \calb=
\bcen
\xymatrix{
\rva\ar[r]\ar[d]
&\rvb\ar[d]
\\
\rvc\ar[r]
\ar[d]
&\rvd
\\
\rve
}
\ecen
\quad
\text{ then }
\left\{
\begin{array}{l}
Fam(\rva) =\{\rva\}
\\
Fam(\rvb)=\{\rvb, \rva\}
\\
Fam(\rvc)=\{\rvc, \rva\}
\\
Fam(\rvd)=\{\rvd, \rvb, \rvc\}
\\
Fam(\rve)=\{\rve, \rvc\}
\\
\calh_{fam}(\calb)=
\{\{\rva\}, \{\rvb, \rva\}, \{\rvc, \rva\}, \{\rvd, \rvb, \rvc\}, \{\rve, \rvc\}\}
\end{array}
\right.
 \eeq
Notice that these hyperedges are precisely the scopes of the conditional probability distributions
$P(x|pa(x))$ in the bnet.

The {\bf primal 2-section of} $\calh_{fam}(\calb)$,
which we will denote by $P2(\calh_{fam}(\calb))$, is constructed by
removing or retaining or replacing from $\calh_{fam}(\calb)$, each
$A\in \calh_{fam}(\calb)$ as follows:
\begin{enumerate}
\item If $|A|<2$, remove
\item If $|A|=2$, keep
\item If $|A|>2$, replace by
a clique.
For example, if $A=\{\rva, \rvb, \rvc\}
\in \calh_{fam}(\calb)$, then
we replace it by the 3 elements
$\{\rva, \rvb\}$, $\{\rvb, \rvc\}$
and $\{\rva, \rvc\}$.
\end{enumerate}
For example

\beq
\text{If}\quad \calb=
\bcen
\xymatrix{
\rva\ar[r]\ar[d]
&\rvb\ar[d]
\\
\rvc\ar[r]\ar[d]
&\rvd
\\
\rve
}
\ecen
\quad
\text{ then }
P2[\calh_{fam}(\calb)]= \{
{ba, ca, db, bc, dc, {\color{red}ec}}
\}
\label{eq-p2-h-fam-b}
\eeq
where we are abbreviating set $\{\rva, \rvb\}$ by $ab$ or $ba$.

For the bnet $\calb$ of
Eq.(\ref{eq-p2-h-fam-b}),
consider the query
$Q=\rvc\perp \rvb\mid \rvd?$.
The scope
$S_Q$ of this query
and its ancestor set $An(S_Q)$
are

\beq
S_Q=\{\rvb, \rvc, \rvd\}
,\quad
An(S_Q)=\{\rva, \rvb, \rvc, \rvd\}
\eeq
Let's denote by $\calb|_{S_Q}$,
the DAG of $\calb$, without probabilities, after we erase
all the nodes not in $S_Q$
and edges to or from them.
We'll call this the {\bf restriction
of bnet $\calb$ to node set $S_Q$}.
Clearly,
if we remove
 from Eq.(\ref{eq-p2-h-fam-b}) the element in red, we get

\beq
P2[\calh_{fam}(\calb|_{S_Q})]=
\{ba, ca, db, bc, dc\}
\eeq
But this is precisely the moralized
graph $AM(\calb, Q)$. Thus

\beq
AM(\calb, Q)= P2[\calh_{fam}(\calb|_{S_Q})]
\eeq